% GENERATED FILE. Edit canonical lessons, then run npm run generate:books.
% canonical-source-hash: fnv1a64:495f2934f6aadfdd
% canonical-lessons: LA-C439-appello, LA-C439-egredior, LA-C439-subduco, LA-C439-deduco, LA-C439-navigabilis

\chapter{To put in, To step out, To haul ashore, To launch, Navigable}
\label{ch:la-to-put-in-to-step-out-to-haul-ashore-to-launch-navigable}
\hlchaptermodality{439}

\begin{chapteropening}
\textbf{By the end of this chapter:} \emph{I can talk about arriving and leaving by ship: putting in, going ashore, hauling a ship up, launching it, and a navigable river.}

Complete the last lesson of chapter 439: I can talk about arriving and leaving by ship: putting in, going ashore, hauling a ship up, launching it, and a navigable river.
\end{chapteropening}

\section[appellō, appellere]{appellō, appellere — to put in (to land), to bring a ship to shore}

\label{lesson:LA-C439-appello}



\begin{quote}
Before the new one: say the Latin for to set sail, then the Latin for a ferryman.
\end{quote}

\subsection*{You'll want to know: appellō, appellere}
\textbf{appellō, appellere} — \textquotedblleft{}to put in (to land), to bring a ship to shore\textquotedblright{}.

\begin{quote}
Nautae nāvem ad lītus appellunt.
\end{quote}

\begin{quote}
\emph{The sailors bring the ship in to the shore.}
\end{quote}

\emph{appellō, appellere, appulī, appulsum}, third conjugation. Do not mix it up with \emph{appellō, appellāre} 'to call', a different verb.

\begin{cousinweb}
\emph{appellō, appellere} is built from \emph{ad-} 'to' and \emph{pellō} 'to drive, to push' (EtymDB; Lewis \& Short): to drive a ship to land.
\end{cousinweb}

\begin{grammarlens}[title={What you've built}]
One more word for telling, explaining and describing at B1.
\end{grammarlens}

\subsection*{Your turn}
\begin{itemize}
\raggedright
  \item \emph{Say it:} \emph{appellō, appellere}
  \item \emph{Say it:} \emph{appellō, appellere}, once more
  \item \emph{Say it:} a sentence of your own with \emph{appellō, appellere}
  \item \emph{Recall:} say the Latin for to set sail, then the Latin for a ferryman, then say \emph{appellō, appellere} again
\end{itemize}

\begin{tcolorbox}[breakable,colback=teal!4,colframe=teal!35!black,title={Before you move on}]
What does appellō, appellere mean? (\textquotedblleft{}To put in (to land), to bring a ship to shore\textquotedblright{}.) Say it once more.
\end{tcolorbox}
\section[ēgredior, ēgredī]{ēgredior, ēgredī — to step out, to go out; to land, to disembark}

\label{lesson:LA-C439-egredior}



\begin{quote}
Before the new one: say the Latin for a ferryman, then the Latin for to put in.
\end{quote}

\subsection*{You'll want to know: ēgredior, ēgredī}
\textbf{ēgredior, ēgredī} — \textquotedblleft{}to step out, to go out; to land, to disembark\textquotedblright{}.

\begin{quote}
Viātōrēs ē nāve ēgrediuntur et ad caupōnam ambulant.
\end{quote}

\begin{quote}
\emph{The travellers get off the ship and walk to the inn.}
\end{quote}

Deponent: \emph{ēgredior, ēgredī, ēgressus sum}. \emph{Ē nāve ēgredī}, to disembark. English \emph{egress} comes from it.

\begin{cousinweb}
\emph{ēgredior} is built from \emph{ē(x)-} 'out' and \emph{gradior} 'to step' (EtymDB; Lewis \& Short). English \emph{egress} comes from its participle \emph{ēgressus}.
\end{cousinweb}

\begin{grammarlens}[title={What you've built}]
One more word for telling, explaining and describing at B1.
\end{grammarlens}

\subsection*{Your turn}
\begin{itemize}
\raggedright
  \item \emph{Say it:} \emph{ēgredior, ēgredī}
  \item \emph{Say it:} \emph{ēgredior, ēgredī}, once more
  \item \emph{Say it:} a sentence of your own with \emph{ēgredior, ēgredī}
  \item \emph{Recall:} say the Latin for a ferryman, then the Latin for to put in, then say \emph{ēgredior, ēgredī} again
\end{itemize}

\begin{tcolorbox}[breakable,colback=teal!4,colframe=teal!35!black,title={Before you move on}]
What does ēgredior, ēgredī mean? (\textquotedblleft{}To step out, to go out; to land, to disembark\textquotedblright{}.) Say it once more.
\end{tcolorbox}
\section[subdūcō, subdūcere]{subdūcō, subdūcere — to haul ashore, to draw up (a ship); to take away}

\label{lesson:LA-C439-subduco}



\begin{quote}
Before the new one: say the Latin for to put in, then the Latin for to step out.
\end{quote}

\subsection*{You'll want to know: subdūcō, subdūcere}
\textbf{subdūcō, subdūcere} — \textquotedblleft{}to haul ashore, to draw up (a ship); to take away\textquotedblright{}.

\begin{quote}
Hieme nautae nāvēs in lītus subdūcunt.
\end{quote}

\begin{quote}
\emph{In winter the sailors haul the ships up on the shore.}
\end{quote}

\emph{subdūcō, subdūcere, subdūxī, subductum}. Ancient ships were pulled up on the beach for the winter.

\begin{cousinweb}
\emph{subdūcō} is built from \emph{sub-} 'from under, up' and \emph{dūcō} 'to lead, to draw' (EtymDB; Lewis \& Short). English \emph{subduction} comes from the noun \emph{subductiō}.
\end{cousinweb}

\begin{grammarlens}[title={What you've built}]
One more word for telling, explaining and describing at B1.
\end{grammarlens}

\subsection*{Your turn}
\begin{itemize}
\raggedright
  \item \emph{Say it:} \emph{subdūcō, subdūcere}
  \item \emph{Say it:} \emph{subdūcō, subdūcere}, once more
  \item \emph{Say it:} a sentence of your own with \emph{subdūcō, subdūcere}
  \item \emph{Recall:} say the Latin for to put in, then the Latin for to step out, then say \emph{subdūcō, subdūcere} again
\end{itemize}

\begin{tcolorbox}[breakable,colback=teal!4,colframe=teal!35!black,title={Before you move on}]
What does subdūcō, subdūcere mean? (\textquotedblleft{}To haul ashore, to draw up (a ship); to take away\textquotedblright{}.) Say it once more.
\end{tcolorbox}
\section[dēdūcō, dēdūcere]{dēdūcō, dēdūcere — to launch (a ship); to lead away, to escort}

\label{lesson:LA-C439-deduco}



\begin{quote}
Before the new one: say the Latin for to step out, then the Latin for to haul ashore.
\end{quote}

\subsection*{You'll want to know: dēdūcō, dēdūcere}
\textbf{dēdūcō, dēdūcere} — \textquotedblleft{}to launch (a ship); to lead away, to escort\textquotedblright{}.

\begin{quote}
Vēre nautae nāvēs in mare dēdūcunt.
\end{quote}

\begin{quote}
\emph{In spring the sailors launch the ships into the sea.}
\end{quote}

\emph{dēdūcō, dēdūcere, dēdūxī, dēductum}. The opposite of \emph{subdūcō}. English \emph{deduce} and \emph{deduct} come from it.

\begin{cousinweb}
\emph{dēdūcō} is built from \emph{dē-} 'down, away' and \emph{dūcō} 'to lead' (EtymDB; Lewis \& Short): to draw a ship down into the water. English \emph{deduce} and \emph{deduct} come from it.
\end{cousinweb}

\begin{grammarlens}[title={What you've built}]
One more word for telling, explaining and describing at B1.
\end{grammarlens}

\subsection*{Your turn}
\begin{itemize}
\raggedright
  \item \emph{Say it:} \emph{dēdūcō, dēdūcere}
  \item \emph{Say it:} \emph{dēdūcō, dēdūcere}, once more
  \item \emph{Say it:} a sentence of your own with \emph{dēdūcō, dēdūcere}
  \item \emph{Recall:} say the Latin for to step out, then the Latin for to haul ashore, then say \emph{dēdūcō, dēdūcere} again
\end{itemize}

\begin{tcolorbox}[breakable,colback=teal!4,colframe=teal!35!black,title={Before you move on}]
What does dēdūcō, dēdūcere mean? (\textquotedblleft{}To launch (a ship); to lead away, to escort\textquotedblright{}.) Say it once more.
\end{tcolorbox}
\section[nāvigābilis, nāvigābile]{nāvigābilis, nāvigābile — navigable}

\label{lesson:LA-C439-navigabilis}



\begin{quote}
Before the new one: say the Latin for to haul ashore, then the Latin for to launch.
\end{quote}

\subsection*{You'll want to know: nāvigābilis, nāvigābile}
\textbf{nāvigābilis, nāvigābile} — \textquotedblleft{}navigable\textquotedblright{}.

\begin{quote}
Flūmen magnum et altum est, itaque nāvigābile.
\end{quote}

\begin{quote}
\emph{The river is big and deep, and so it is navigable.}
\end{quote}

An adjective of two endings, from \emph{nāvigō} 'to sail' with \emph{-bilis} 'that can be'. English \emph{navigable} comes from it.

\begin{cousinweb}
\emph{nāvigābilis} is built on \emph{nāvigō} 'to sail' with the suffix \emph{-bilis} 'that can be' (Lewis \& Short). English \emph{navigable} comes from it.
\end{cousinweb}

\begin{grammarlens}[title={What you've built}]
One more word for telling, explaining and describing at B1.
\end{grammarlens}

\subsection*{Your turn}
\begin{itemize}
\raggedright
  \item \emph{Say it:} \emph{nāvigābilis, nāvigābile}
  \item \emph{Say it:} \emph{nāvigābilis, nāvigābile}, once more
  \item \emph{Say it:} a sentence of your own with \emph{nāvigābilis, nāvigābile}
  \item \emph{Recall:} say the Latin for to haul ashore, then the Latin for to launch, then say \emph{nāvigābilis, nāvigābile} again
\end{itemize}

\begin{tcolorbox}[breakable,colback=teal!4,colframe=teal!35!black,title={Before you move on}]
What does nāvigābilis, nāvigābile mean? (\textquotedblleft{}Navigable\textquotedblright{}.) Say it once more.
\end{tcolorbox}
